\begin{abstract}
Conditional flow matching (CFM) turns generative modelling into a least-squares regression, but each sample costs one network evaluation per ODE step (NFE).
With one network and one training budget fixed, we ask on MNIST which changes reduce the NFE most per unit of training cost.
Training seeds move FID by up to 7.8 points (sd) against under 0.6 for sampling noise, and a $k$-means class-coverage test with a two-sample $\chi^2$ test shows that FID and class balance can disagree.
A time grid derived from the local Euler error (step density $\propto\sqrt{\kappa}$) and second-order solvers whose last step avoids the network at $t=1$ cut the NFE needed for FID 20 by 28--58\% at no training cost.
Minibatch optimal transport saves at most 6\% of the pairing cost (a covariance bound we derive allows 26\%), and larger OT batches do not improve FID; logit-normal training times match the best sampler's gain (one training seed).
Reflow makes the paths almost straight and lowers one-step FID from 332 to 65 (I-CFM) and from 223 to 61 (OT-CFM) for 2.7--2.8 times the training time, but raises the many-step FID to about 21.
Reflow is the best choice for one to four steps, improved samplers or training times from about seven steps on.
\end{abstract}
